\documentclass[11pt]{article}
\usepackage[margin=1in]{geometry}
\usepackage{amsmath,amssymb,amsthm}
\usepackage{bm}
\usepackage{booktabs}
\usepackage{hyperref}
\hypersetup{colorlinks=true,linkcolor=blue,urlcolor=blue}

\newcommand{\R}{\mathbf{R}}
\newcommand{\D}{\mathbf{D}}
\newcommand{\bet}{\bm{\beta}}
\newcommand{\Ai}{A_i}
\newcommand{\E}{\mathbb{E}}
\newcommand{\Var}{\operatorname{Var}}
\newcommand{\Cov}{\operatorname{Cov}}

\title{\textbf{Spatially continuous (raster) covariates by intensity-scale
aggregation}\\[4pt]
\large Why averaging predictors over polygons is biased, and the
\texttt{SDALGCP2\_raster} fix}
\author{SDALGCP2 development notes}
\date{\today}

\begin{document}
\maketitle

\begin{abstract}
When a covariate varies \emph{within} an areal unit it is usually summarised by its
areal mean and entered into a Poisson model. Under the nonlinear log link this is
biased. We derive the statistically coherent alternative used by
\texttt{SDALGCP2\_raster}: covariates supplied as rasters enter the log-Gaussian
Cox process intensity at the candidate-point level and are aggregated on the
\emph{intensity} (exponential) scale through a log-sum-exp offset. We give the
exact discretised model, the factorisation that isolates the covariate
contribution, the Gauss--Newton fixed point used to fit it (which reuses the
ordinary SDA-LGCP machinery), and a simulation in which areal averaging is biased
by $+67\%$ while the intensity-scale fit recovers the truth.
\end{abstract}

\section{The aggregated LGCP mean}

Let the LGCP intensity be $\lambda(x)=\rho(x)\exp\{z(x)^\top\bet+S(x)\}$, with
$\rho(x)$ the population (at-risk) density, $z(x)$ a vector of covariates available
as rasters, and $S$ a zero-mean Gaussian process with
$\Cov\{S(x),S(x')\}=\sigma^2\rho_\phi(\|x-x'\|)$. The count in region $\Ai$ is
conditionally Poisson with mean $\int_{\Ai}\lambda(x)\,dx$. Writing
$m_i=\int_{\Ai}\rho(x)\,dx$ for the population at risk and
$w_i(x)=\rho(x)/m_i$ (so $\int_{\Ai}w_i=1$),
\begin{equation}
\E[Y_i\mid S]=m_i\int_{\Ai} w_i(x)\,\exp\!\big\{z(x)^\top\bet+S(x)\big\}\,dx .
\label{eq:mean}
\end{equation}
With the same candidate points $x_{ik}$ and quadrature weights $w_{ik}$ used to
build the aggregated correlation (companion note on the scale parameter),
\begin{equation}
\E[Y_i\mid S]\ \approx\ m_i\sum_{k=1}^{n_i} w_{ik}\,
   \exp\!\big\{z(x_{ik})^\top\bet + S(x_{ik})\big\}.
\label{eq:disc}
\end{equation}

\section{Why averaging the predictor over polygons is biased}

The conventional ``areal'' model replaces $z(x)$ by its within-region mean
$\bar z_i=\sum_k w_{ik}z(x_{ik})$ and sets the region linear predictor to
$\bar z_i^\top\bet$. But by Jensen's inequality, for the covariate part alone,
\begin{equation}
\sum_k w_{ik}\exp\{z(x_{ik})^\top\bet\}
\ \ge\ \exp\{\bar z_i^\top\bet\},
\end{equation}
with equality only if $z$ is constant on $\Ai$. The gap is governed by the
\emph{within-region} variance of $z^\top\bet$. Crucially, a regression of $Y_i$ on
$\bar z_i$ is biased \emph{for the coefficient} $\bet$ precisely when that
within-region variance is correlated with $\bar z_i$ across regions --- for example
when the covariate has sharp localised peaks (point exposure sources): regions
containing a peak have a moderate mean but a large intensity aggregate. Averaging
washes the peaks out; the exponential does not.

\section{Intensity-scale factorisation}

Factor the covariate out of \eqref{eq:disc}:
\begin{equation}
\boxed{\;
\E[Y_i\mid S]= m_i\,e^{\,b_i(\bet)}\sum_k c_{ik}(\bet)\,e^{S(x_{ik})},
\qquad
b_i(\bet)=\log\!\sum_k w_{ik}\,e^{\,z(x_{ik})^\top\bet}\;}
\label{eq:factor}
\end{equation}
where the \emph{intensity-tilted} weights
\begin{equation}
c_{ik}(\bet)=\frac{w_{ik}\,e^{\,z(x_{ik})^\top\bet}}{\sum_l w_{il}\,e^{\,z(x_{il})^\top\bet}},
\qquad \sum_k c_{ik}=1,
\label{eq:tilt}
\end{equation}
upweight high-intensity sub-areas. The term $b_i(\bet)$ is the \textbf{log-sum-exp}
of the within-region covariate surface; it equals $\bar z_i^\top\bet$ only when $z$
is constant on $\Ai$, and is the correct region-level covariate contribution.

\section{Reducing the spatial part}

Let $T_i=\sum_k c_{ik}S(x_{ik})$, a Gaussian linear functional, and write
$S(x_{ik})=T_i+\delta_{ik}$ with $\sum_k c_{ik}\delta_{ik}=0$. A second-order
expansion gives
\begin{equation}
\sum_k c_{ik}e^{S(x_{ik})}=e^{T_i}\sum_k c_{ik}e^{\delta_{ik}}
\approx e^{T_i}\Big(1+\tfrac12\sum_k c_{ik}\delta_{ik}^2\Big)
\approx \exp\!\big(T_i+\tfrac12 v_i\big),
\end{equation}
with $v_i=\E\sum_k c_{ik}\delta_{ik}^2=\sigma^2\big(1-\sum_{k,l}c_{ik}c_{il}\rho_\phi(d_{kl})\big)
=\sigma^2\big(1-\R^c_{ii}\big)$, the variance lost to spatial averaging. Hence the
\emph{full} (covariate-tilted) model is
\begin{equation}
\E[Y_i\mid T_i]= m_i\exp\!\big(b_i(\bet)+\tfrac12 v_i(\bet)+T_i\big),
\quad
\mathbf{T}\sim\mathcal N\!\big(\mathbf 0,\ \sigma^2\R^c(\phi,\bet)\big),
\quad
\R^c_{ij}=\sum_{k,l}c_{ik}c_{jl}\rho_\phi(d_{kl}).
\label{eq:full}
\end{equation}

\paragraph{Offset-only model (the default implementation).}
Because the dominant correction is the covariate term $b_i$, the implemented
default keeps the spatial aggregation at the population weights
($\R^c\!\to\!\R^w$, $\tfrac12 v_i$ absorbed), so the correlation matrix does
\emph{not} depend on $\bet$ and is precomputed once:
\begin{equation}
\E[Y_i\mid T_i]= m_i\exp\!\big(b_i(\bet)+T_i\big),
\qquad
\mathbf T\sim\mathcal N\!\big(\mathbf 0,\ \sigma^2\R^w(\phi)\big).
\label{eq:offset}
\end{equation}
This captures the covariate Jensen bias \emph{exactly} through $b_i(\bet)$ and
treats the covariate tilting of the latent field as second order. The fully tilted
model \eqref{eq:full} is available as a refinement (it rebuilds $\R^c(\bet)$ each
outer iteration).

\section{Estimation: a Gauss--Newton fixed point}

The only nonlinearity in $\bet$ is through $b_i(\bet)$, whose gradient is the
\emph{intensity-tilted (``effective'') covariate}
\begin{equation}
\frac{\partial b_i}{\partial \bet}
=\frac{\sum_k w_{ik}e^{z_{ik}^\top\bet}z_{ik}}{\sum_k w_{ik}e^{z_{ik}^\top\bet}}
=\sum_k c_{ik}(\bet)\,z(x_{ik})\ \equiv\ \D^c_i(\bet).
\label{eq:effdesign}
\end{equation}
At iteration $t$ we linearise $b_i(\bet)\approx b_i(\bet^{(t)})
+\D^c_i(\bet^{(t)})^\top(\bet-\bet^{(t)})$ and fit the ordinary SDA-LGCP MCML with
\emph{effective design} $\D^c(\bet^{(t)})$ and offset
$o_i^{(t)}=\log m_i+b_i(\bet^{(t)})-\D^c_i(\bet^{(t)})^\top\bet^{(t)}$, obtaining
$\bet^{(t+1)}$ (and $\sigma^2,\phi$). At a fixed point $\bet^\star$ the offset
reproduces $b_i(\bet^\star)$ exactly and, by \eqref{eq:effdesign}, the linear score
$\sum_i\D^c_i(\cdot)$ equals the nonlinear score $\sum_i(\partial b_i/\partial\bet)(\cdot)$,
so $\bet^\star$ solves the likelihood equations of the nonlinear model. The scheme
therefore reuses \texttt{mcml\_fit()} unchanged and is exactly Gauss--Newton on the
$b_i(\bet)$ term. (\texttt{SDALGCP2\_raster}; the iteration converges in a handful
of steps.)

\section{Numerical demonstration}

A covariate with $14$ sharp Gaussian peaks (width $\sim0.5$, region width $2.5$) on
an $8\times8$ lattice; counts simulated from \eqref{eq:disc} with $\beta_z=1$.

\begin{center}
\begin{tabular}{@{}lcc@{}}
\toprule
method & $\widehat{\beta_z}$ & bias \\
\midrule
truth & $1.00$ & --- \\
naive (areal-mean covariate, Poisson GLM) & $1.67$ & $+67\%$ \\
exact nonlinear MLE of \eqref{eq:disc} & $1.01$ & $+1\%$ \\
\texttt{SDALGCP2\_raster} (offset model) & $1.06$ & $+6\%$ \\
\bottomrule
\end{tabular}
\end{center}

Areal averaging is badly biased; intensity-scale aggregation recovers the truth and
matches the exact MLE. The bias is large only when the within-region covariate
variance correlates with the region mean (sharp sub-region features); for smooth
covariates the two approaches nearly agree, as expected from
$b_i(\bet)\approx\bar z_i^\top\bet+\tfrac12\bet^\top\widehat{\Var}_i(z)\bet$.

\paragraph{Reference.}
Johnson, O., Diggle, P., Giorgi, E. (2019). A spatially discrete approximation to
log-Gaussian Cox processes for aggregated disease count data. \emph{Statistics in
Medicine} 38, 4871--4887.

\end{document}
